\section{Invariant differentials}

Let $F$ be a formal group law over $R$. In this setting, a differential form is an expression
of the form
$P(X) \, dX$ for some $P(X) \in R[\![X]\!]$. In this section, we discuss those
differential forms that respect the group structure of $F$.

\begin{definition}[{\cite[Chap.~4 Definition 4.1]{silverman2009arithmetic}}]
    An \emph{invariant differential} on a formal group law $F$ over a ring $R$ is a differential form
    $\omega (X) = P(X)\, dX$ with $P(X) \in R[\![X]\!]$ satisfying
    \[\omega(X) \circ F(X,Y) = \omega (X).\]
    To be precise, $P(F(X,Y)) F_X (X,Y) = P(X)$, where $F_X(X,Y)$ is the partial derivative of
    $F$ with respect to its first variable. An invariant differential is called \emph{normalized}
    if $P(0) = 1$.
\end{definition}

Recall that a power series $f$ over $R$ is invertible in $R[\![X]\!]$
if and only if its constant coefficient is invertible in $R$.

\begin{proposition}[{\cite[Chap.~4 Proposition 4.2]{silverman2009arithmetic}}]
    Let $F$ be a formal group law over $R$. Then there exists a unique normalized invariant
    differential on $F$. It is given by the formula
    \[ \omega (X) = F_X (0,X) ^ {-1} \, dX. \]
    Moreover, every invariant differential on $F$ can be written as $a \omega$ for some $a \in R$.
\end{proposition}

\begin{proof}
    Write $\omega (X) = P(X)\, dX$. Since $\omega$ is invariant, we must have
    \[ P(F(X,Y)) F_X(X,Y) = P (X). \]
    Let $X = 0$ in the above equation and use the fact that $F(0 ,Y) = Y$.
    Then we have
    \[ P(Y) F_X(0, Y) = P (0). \]
    Since $F_X(0,Y)$ is of the form $(1 + O(Y))$, $F_X(0,Y)$ is invertible in
    $R[\![Y]\!]$. Here $O(Y)$ is a power series of order at least $1$. Then
    \begin{equation}
        P(Y) = P(0) F_X(0,Y) ^ {-1}.
        \label{eq: diff_eq}
    \end{equation}
    i.e. $P(X)$ is completely determined by the value $P(0)$.
    Therefore, it is enough to verify that
    \[ \omega = F_X(0,X) ^ {-1}\, dX. \]
    is a normalized invariant differential.

    It is normalized since $F_X(0,X)$ is of the form $(1 + O(X))$. Hence, we just need to
    show that $\omega$ is invariant, i.e.
    \[ F_X(0, F(Y,Z)) ^{-1} F_X(Y,Z) = F_X(0,Y) ^ {-1}. \]
    Consider the associativity condition $F(F(X,Y),Z) = F(X,F(Y,Z))$ and differentiate it
    with respect to $X$. Then we have
    \[F_X(X,F(Y,Z)) = F_X(F(X,Y), Z) F_X(X,Y). \]
    Taking $X = 0$ above, we have
    \[F_X(0,F(Y,Z)) = F_X(Y,Z) F_X(0,Y). \]
    This is exactly what we want.

    Finally, by equation \eqref{eq: diff_eq}, it follows that every invariant differential is
    of form $a \omega$.
\end{proof}

\begin{example}
    The normalized invariant differential on $\GG_a$ is $dX$, and the normalized invariant
    differential on $\GG_m$ is $(1 + X) ^{-1}\, dX = \sum_{n=0}^{\infty} (-1)^n X ^ n\, d X$.
\end{example}

\begin{corollary}[{\cite[Chap.~4 Corollary 4.3]{silverman2009arithmetic}}]
    Let $F,G$ be formal group laws over $R$ with normalized invariant differentials $\omega_F$
    and $\omega_G$, respectively. Let $f$ be a homomorphism from $F$ to $G$. Then
    \[ \omega_G \circ f = f'(0) \omega_F .\]
    where $f'(X)$ is the formal derivative of $f(X)$.
    \label{cor: omega_hom_a}
\end{corollary}

\begin{proof}
    According to the definition, we have
    \[ \omega_G \circ f (F (X,Y)) = \omega_G \circ G(f(X), f(Y)) = \omega_G \circ f (X).\]
    Then $\omega_G \circ f$ is an invariant differential on $F$. It follows that $\omega_G \circ f$
    can be expressed as $a \omega_F$ for some $a \in R$. By checking the coefficient of $X$,
    we can deduce that $a = f'(0)$.
\end{proof}

\begin{corollary}[{\cite[Chap.~4 Corollary 4.4]{silverman2009arithmetic}}]
    Let $F$ be a formal group law over $R$ and let $p$ be a prime number. Then there are power series
    $f(X), g(X)$ over $R$ with zero constant coefficient such that
    \[ [p]_F (X) = pf(X) + g(X^p). \]

    \label{cor: p_series_form}
\end{corollary}

\begin{proof}
    Let $\omega$ be the normalized invariant differential on $F$. We have discussed before
    that $[p](X) = pX + O(X^2)$. Then by Corollary \ref{cor: omega_hom_a},
    together with $[p]'(0) = p$, we have
    \[ \omega \circ [p] = p \omega. \]
    If we write $\omega(X) = P(X) \, dX$, then
    \[\omega \circ [p] (X) = P(X) [p]'(X) \, dX. \]
    It follows that $P(X) [p]'(X) \, dX = p \omega(X) = p P(X)\, dX. $
    Since $\omega$ is a normalized invariant differential on $F$, then $P(X)$ is of the
    form $(1 + O(X))$. It follows that $P(X)$ is invertible in
    $R[\![X]\!]$. Therefore, we have
    \[[p]'(X) \in p R[\![X]\!]. \]
    If we denote $[p](X) = \sum_{n=1}^{\infty} a_n X ^ n$. Then for each $n \in \NN_{>0}$
    such that $p \nmid n$, then $a_n \in p R$.
    % $a_n$ satisfies either $p \mid n$ or $a_n \in p R$.
\end{proof}

\begin{example}
    Assume that $R$ is a ring of characteristic $p$. Then the $p$-series $[p]$ has the form
    $[p](X) = g(X ^ p)$ for some $g(X) \in R [\![X]\!]$.
\end{example}